\section{Related Work}
\label{sec:related-work}

\textsc{QuaRK} sits at the intersection of reservoir computing, quantum
feature generation, kernel learning, and statistical learning from dependent
data. The method does not replace the central ideas from these areas; rather,
it combines them in a temporal architecture whose dynamical, statistical, and
measurement assumptions can be tracked jointly.

\paragraph{Reservoir computing, causal filters, and fading memory.}
Echo-state networks and liquid-state machines established the principle of
driving a fixed recurrent system and training only its readout
\citep{Jaeger2001,MaassNatschlagerMarkram2002}. Fading memory formalizes
continuity of a causal filter with respect to increasingly remote inputs
\citep{BoydChua1985}, while universality and risk analyses study, respectively,
the functions representable by reservoir families and the statistical
behavior of learned readouts
\citep{GrigoryevaOrtega2018,GononGrigoryevaOrtega2020}. In the present work,
uniform contraction is used for the narrower purpose of proving loss of
initialization dependence and uniqueness of the causal state for one
reset-based architecture. We do not rely on a separate universality or formal
fading-memory theorem.

\paragraph{Quantum reservoir computing.}
Quantum reservoir computing realizes the driven state system with quantum
dynamics \citep{FujiiNakajima2017}. Existing proposals have considered
dissipative evolution, reset mechanisms, noisy processors, weak or repeated
measurements, and spatial multiplexing
\citep{ChenNurdin2019,ChenNurdinYamamoto2020,NakajimaEtAl2019,
MolteniDestriPrati2023,MujalEtAl2023,HuEtAl2024}. Finite-dimensional quantum
reservoir filters and the role of input dependence have also received formal
treatment \citep{MartinezPenaOrtega2023}. \textsc{QuaRK} uses a replacement
channel whose contraction coefficient is exact and tunable. Its multiple
reservoir instances are modeled as separate state evolutions, so they may be
executed in parallel or sequentially on reused hardware; this distinction is
important for separating peak width from runtime and state-preparation cost.

\paragraph{Quantum features and kernel readouts.}
Quantum feature methods encode classical data into quantum states or
measurement statistics before applying a classical learner
\citep{SchuldKilloran2019,HuangBroughtonEtAl2021}. Here the features are
history-dependent: they are generated after a quantum reservoir has processed
a temporal window. The downstream Mat\'ern kernel is an ordinary classical
positive-definite kernel on these measured features. Kernel ridge regression
provides a regularized nonlinear readout, while the normalized RKHS structure
bounds predictions and allows feature perturbations to be propagated to the
readout \citep{Genton2001,PorcuEtAl2024,HofmannScholkopfSmola2008}. Thus the
quantum component lies in the temporal featurizer, not in a claim that the
Mat\'ern kernel lacks a classical analogue.

\paragraph{Projection and measurement.}
The JL lemma preserves pairwise geometry on a finite set under a random linear
map whose target dimension grows logarithmically with the set size
\citep{DasguptaGupta2003,Woodruff2014}. We use it to relate the number of
encoded coordinates to the finite collection of points being processed; the
classical projection cost and the subsequent nonlinear dynamics remain part of
the method. On the measurement side, classical shadows estimate many
observables from shared randomized measurements \citep{HuangKuengPreskill2020}.
For local random Pauli measurements, a Pauli string of weight at most \(k\) has
squared shadow norm scaling as \(3^k\), which determines the locality dependence
in our shot bound.

\paragraph{Learning with dependent observations.}
Independent-block methods compare a beta-mixing sequence with a collection of
independent blocks and yield uniform learning bounds in terms of an effective
block count \citep{Yu1994,MohriRostamizadeh2008,DedeckerEtAl2007}. We apply this
machinery to strided temporal windows. The resulting confidence condition
depends on the raw-time separation between blocks, and the stochastic terms
scale with the number of approximately independent blocks rather than the raw
number of windows.

The standard ingredients above become useful for \textsc{QuaRK} only after
they are connected to the same temporal feature map. The technical
contribution is therefore their architecture-specific integration: tuple-based
multiplexing, finite-window error propagation through the Mat\'ern readout, a
window-process mixing calculation that retains the effective block count, and
separate accounting of exact-feature learning and finite-shot measurement. We
now introduce the stochastic objects needed to state this connection
precisely.
